\documentclass[11pt,a4paper]{article}  % Déclare la classe du document.

\newcommand{\chaptermark}[2]{\og #2 \fg{} (#1)} % ligne uniquement pour éviter une erreur en cas d'article
\include{../1ere-preambule}

\usepackage{epic}   % extension et simplification de  l'environnement picture
\usepackage{graphicx}    % permet d'inclure des graphique externe dans un document
\usepackage{arydshln}   %permet d'avoir des lignes en pointillés

\newcommand{\lignes}[1]{
	\setlength{\unitlength}{1mm}
	\vskip 0,4cm\noindent
	\multido{\i=0+1}{#1}{%
		\noindent
		\begin{picture}(150,5)
			\linethickness{0,005mm}
			\put(-5,0){\dashline{1}(175,5)(5,5)}
		\end{picture}\vskip 0,001cm
	}
	\vskip 0cm
}

\rhead{\red{Correction - DS n$^{\circ}02$ - Trigonométrie - Sujet B} - 28/11/2022}

\newcommand{\va}[1]{\lvert \ #1 \ \rvert}
\newcommand{\vaf}[1]{\left\lvert \ #1\  \right\rvert}

\begin{document}
	\graphicspath{{images/}}
	
	\begin{center}
		
		
		\fbox{
			\begin{minipage}[]{15cm}
				\begin{center}
					\LARGE{\rule{0.5cm}{0pt}\rule[-0.25cm]{0pt}{0.9cm}
						\red{Correction - DS n$^{\circ}02$} \\
						Trigonométrie - Sujet B 
					}
				\end{center}
			\end{minipage}
		}
	\end{center}

\begin{center}\textbf{1h50min - Barème indicatif}\end{center}
\begin{center}\textbf{Les élèves avec un tier-temps ne traitent pas les questions avec  le symbole \ding{46}}\end{center}
\begin{center}\textbf{Calculatrice autorisée.\\Les résultats devront être justifiés (à l'aide de calcul ou de propriétés).}\end{center}


\exerciceDS {- enroulement autour du cercle trigo - 4 points }

En utilisant la figure ci-dessous, donner les points du cercle qui correspondent aux angles suivants. On détaillera la décomposition du nombre en donnant la mesure principale de l'angle.
%\begin{multicols}{4}
	\begin{enumMet}
		\item $\dfrac{439\pi}{6}$\quad \corr{$=\dfrac{444\pi}{6}-\dfrac{5\pi}{6}=37\times2\pi+$\fbox{$-\dfrac{5\pi}{6}$}. La mesure principale est $-\dfrac{5\pi}{6}$. Il s'agit du point G.}
		\item $\dfrac{197\pi}{4}$\quad \corr{$=\dfrac{200\pi}{4}-\dfrac{3\pi}{4}=25\times2\pi+$\fbox{$\dfrac{-3\pi}{4}$}. La mesure principale est $\dfrac{-3\pi}{4}$. Il s'agit du point H.}
		\item $\dfrac{-73\pi}{3}$\quad \corr{$=-\dfrac{72\pi}{3}-\dfrac{\pi}{3}=-12\times2\pi+$\fbox{$\dfrac{-\pi}{3}$}. La mesure principale est $\dfrac{-\pi}{3}$. Il s'agit du point L.}
		\item \ding{46} $-\dfrac{617\pi}{2}$\quad \corr{$=-\dfrac{616\pi}{2}-\dfrac{\pi}{2}=-154\times2\pi+$\fbox{$\dfrac{-\pi}{2}$}. La mesure principale est $\dfrac{-\pi}{2}$. Il s'agit du point J'.}
	\end{enumMet}
%\end{multicols}
\ \\
\begin{center}
	\newrgbcolor{ududff}{0.30196078431372547 0.30196078431372547 1.}
	\newrgbcolor{qqwuqq}{0. 0.39215686274509803 0.}
	\psset{xunit=4.0cm,yunit=4.0cm,algebraic=true,dimen=middle,dotstyle=o,dotsize=5pt 0,linewidth=2.pt,arrowsize=3pt 2,arrowinset=0.25}
	\begin{pspicture*}(-1.1,-1.1550625290940115)(1.1008598498816817,1.2224219662799085)
		\psaxes[labelFontSize=\scriptstyle,xAxis=true,yAxis=true,Dx=1.,Dy=1.,ticksize=-2pt 0,subticks=2]{->}(0,0)(-1.98066,-1.1550625290940115)(1.1008598498816817,1.2224219662799085)
		\pscircle[linewidth=2.pt](0.,0.){4.}
		\psline[linewidth=2.pt,linestyle=dashed,dash=1pt 1pt,linecolor=red](0.,0.)(0.8660254037844387,0.5)
		\psline[linewidth=2.pt,linestyle=dashed,dash=1pt 1pt,linecolor=red](0.,0.)(-0.8660254037844387,0.5)
		\psline[linewidth=2.pt,linestyle=dashed,dash=1pt 1pt,linecolor=red](0.,0.)(-0.8660254037844386,-0.5)
		\psline[linewidth=2.pt,linestyle=dashed,dash=1pt 1pt,linecolor=red](0.,0.)(0.8660254037844384,-0.5)
		\psline[linewidth=2.pt,linestyle=dashed,dash=1pt 1pt,linecolor=blue](0.,0.)(0.7071067811865476,0.7071067811865475)
		\psline[linewidth=2.pt,linestyle=dashed,dash=1pt 1pt,linecolor=blue](0.,0.)(-0.7071067811865475,0.7071067811865476)
		\psline[linewidth=2.pt,linestyle=dashed,dash=1pt 1pt,linecolor=blue](0.,0.)(-0.7071067811865477,-0.7071067811865475)
		\psline[linewidth=2.pt,linestyle=dashed,dash=1pt 1pt,linecolor=blue](0.,0.)(0.7071067811865474,-0.7071067811865477)
		\psline[linewidth=2.pt,linestyle=dotted,linecolor=qqwuqq](0.,0.)(0.5,0.8660254037844386)
		\psline[linewidth=2.pt,linestyle=dotted,linecolor=qqwuqq](0.,0.)(-0.5,0.8660254037844387)
		\psline[linewidth=2.pt,linestyle=dotted,linecolor=qqwuqq](0.,0.)(-0.5,-0.8660254037844385)
		\psline[linewidth=2.pt,linestyle=dotted,linecolor=qqwuqq](0.,0.)(0.5,-0.8660254037844386)
		\begin{scriptsize}
			\psdots[dotsize=4pt 0,dotstyle=*,linecolor=darkgray](0.,0.)
			\rput[bl](0.00960106297279085,0.042848354149331562){\Large\black{$O$}}
			\psdots[dotstyle=*,linecolor=ududff](1.,0.)
			\rput[bl](1.027748218996967,0.03743267192752309){\Large\black{$I$}}
			\psdots[dotstyle=*,linecolor=ududff](0.8660254037844387,0.5)
			\rput[bl](0.8761092808657066,0.5265068007775197){\Large\black{$A$}}
			\psdots[dotstyle=*,linecolor=ududff](0.7071067811865476,0.7071067811865475)
			\rput[bl](0.7190546663726156,0.7350105663171459){\Large\black{$B$}}
			\psdots[dotstyle=*,linecolor=ududff](0.5,0.8660254037844386)
			\rput[bl](0.5105511264421327,0.8920653507495917){\Large\black{$C$}}
			\psdots[dotstyle=*,linecolor=ududff](0.,1.)
			\rput[bl](0.00960106297279085,1.0409966118493248){\Large\black{$J$}}
			\psdots[dotstyle=*,linecolor=ududff](-0.5,0.8660254037844387)
			\rput[bl](-0.5427979259339429,0.9028967151932086){\Large\black{$D$}}
			\psdots[dotstyle=*,linecolor=ududff](-0.7071067811865475,0.7071067811865476)
			\rput[bl](-0.7658857895025951,0.7377184074280502){\Large\black{$E$}}
			\psdots[dotstyle=*,linecolor=ududff](-0.8660254037844387,0.5)
			\rput[bl](-0.9256482421766014,0.5265068007775197){\Large\black{$F$}}
			\psdots[dotstyle=*,linecolor=ududff](-1.,0.)
			\rput[bl](-1.0994558275842002,0.020140513038427326){\Large\black{$I'$}}
			\psdots[dotstyle=*,linecolor=ududff](-0.8660254037844386,-0.5)
			\rput[bl](-1.0000187858048398,-0.5597649802551862){\Large\black{$G$}}
			\psdots[dotstyle=*,linecolor=ududff](-0.7071067811865477,-0.7071067811865475)
			\rput[bl](-0.850548494949918,-0.8072236335711421){\Large\black{$H$}}
			\psdots[dotstyle=*,linecolor=ududff](-0.5,-0.8660254037844385)
			\rput[bl](-0.620921440476689,-0.9859411468908217){\Large\black{$K$}}
			\psdots[dotstyle=*,linecolor=ududff](0.,-1.)
			\rput[bl](0.036679444781944466,-1.1094567257687464){\Large\black{$J'$}}
			\psdots[dotstyle=*,linecolor=ududff](0.5,-0.8660254037844386)
			\rput[bl](0.5376295082512864,-0.9205254646690133){\Large\black{$L$}}
			\psdots[dotstyle=*,linecolor=ududff](0.7071067811865474,-0.7071067811865477)
			\rput[bl](0.7488408863626846,-0.7445157924602379){\Large\black{$M$}}
			\psdots[dotstyle=*,linecolor=ududff](0.8660254037844384,-0.5)
			\rput[bl](0.9004798244939449,-0.51305963446988031){\Large\black{$N$}}
		\end{scriptsize}
	\end{pspicture*}
\end{center}
\vspace*{-0.3cm}
\exerciceDS {- même point image - 3 points}
Pour chacun des nombres suivants, déterminer deux autres réels ayant le même point image lors de l'enroulement de la droite numérique
%\begin{multicols}{3}
	\begin{enumMet}
		\item $-\dfrac{4\pi}{7}$\quad  \corr{$-\dfrac{4\pi}{7}+2\pi=-\dfrac{4\pi}{7}+\dfrac{14\pi}{7}=$ \fbox{$\dfrac{10\pi}{7}$} ; $-\dfrac{4\pi}{7}-2\pi=-\dfrac{4\pi}{7}-\dfrac{14\pi}{7}=$ \fbox{$-\dfrac{18\pi}{7}$}}
		\item $\dfrac{3\pi}{5}$\quad  \corr{$\dfrac{3\pi}{5}+2\pi=\dfrac{3\pi}{5}+\dfrac{10\pi}{5}=$ \fbox{$\dfrac{13\pi}{5}$} ; $\dfrac{3\pi}{5}-2\pi=\dfrac{3\pi}{5}-\dfrac{10\pi}{5}=$ \fbox{$-\dfrac{7\pi}{5}$}}
		\item $\dfrac{5\pi}{9}$\quad  \corr{$\dfrac{5\pi}{9}+2\pi=\dfrac{5\pi}{9}+\dfrac{18\pi}{9}=$ \fbox{$\dfrac{23\pi}{9}$} ; $\dfrac{5\pi}{9}-2\pi=\dfrac{5\pi}{9}-\dfrac{18\pi}{9}=$ \fbox{$-\dfrac{13\pi}{9}$}}
	\end{enumMet}
%\end{multicols}

\exerciceDS {- convertir en degrés - 1 point }
Pour chacun des angles suivants, convertir l'angle en degré. On arrondira si nécessaire à 1 chiffre après la virgule
\begin{multicols}{2}
	\begin{enumMet}
		\item $\dfrac{17\pi}{20}$ rad \quad \corr{$=\dfrac{17\pi}{20}\times \dfrac{180}{\pi}=153\degree$}
		\item $\dfrac{7\pi}{13}$ rad \quad \corr{$=\dfrac{7\pi}{13}\times \dfrac{180}{\pi}=96,92\degree$}
	\end{enumMet}
\end{multicols}

\exerciceDS {- convertir en radians - 1 point }
Pour chacun des angles suivants, convertir l'angle en radian. On donnera la valeur exacte/
\begin{multicols}{2}
	\begin{enumMet}
		\item $42\degree$\quad \corr{$=42\times \dfrac{\pi}{180}=\dfrac{7\pi}{30}$ rad}
		\item $523\degree$\quad \corr{$=523\times \dfrac{\pi}{180}=\dfrac{523\pi}{180}$ rad}
	\end{enumMet}
\end{multicols}

\exerciceDS {- cosinus et sinus - 4 points}

En utilisant le cercle trigonométrique de l'exercice 1, recopier et compléter le tableau suivant:\\
\begin{center}
	\begin{tabular}{ | Xc | Xc | Xc | Xc | Xc | }
		\hline
		$x$  & $\dfrac{31\pi}{4}$ & $\dfrac{4\pi}{3}$ & $-\dfrac{25\pi}{6}$ &  $\dfrac{67\pi}{2}$ \\ \hline
		Point image  & \red{M} & \red{K}  & \red{N}  &  \red{J} \\ \hline
		$\cos(x)$  & \red{$\dfrac{\sqrt{2}}{2}$}  & \red{$-\dfrac{1}{2}$}  & \red{$\dfrac{\sqrt{3}}{2}$}  & \red{0}  \\ \hline
		$\sin(x)$ & \red{$-\dfrac{\sqrt{2}}{2}$}  & \red{$-\dfrac{\sqrt{3}}{2}$}  & \red{$-\dfrac{1}{2}$}  & \red{1}  \\
		\hline  
	\end{tabular}
\end{center}

\exerciceDS {- équations trigonométriques - 7 points}

Résoudre les équations suivantes : 
\begin{enumMet}
	\item \ding{46} $\cos(x)=\dfrac{\sqrt{3}}{2}$ avec $x\in \left[0;2\pi\right]$\\
	\begin{corrBox}
		On commence par étudier la solution dans l'intervalle $]-\pi;\pi]$ (parmi les mesures principales de l'angle).\\
		L'équation $\cos(x)=\dfrac{\sqrt{3}}{2}$ a pour solution :\\
		$\begin{cases}
			x_1 &=-\dfrac{\pi}{6}\\
			x_2 &=\dfrac{\pi}{6}
		\end{cases}$\\
		On en déduit les solutions sur \R.\\
		L'équation $\sin(x)=\dfrac{\sqrt{3}}{2}$ a pour solution ($k\in \Z$ ) :\\
		$\begin{cases}
			x_1 &=-\dfrac{\pi}{6}+2k\pi\\
			x_2 &=\dfrac{\pi}{6}+2k\pi\
		\end{cases}$\\
		On teste différentes valeurs de $k$ pour trouver les valeurs qui appartiennent à l'intervalle souhaité $\left[0;2\pi\right]$ (ou $\left[\dfrac{0\pi}{6};\dfrac{12\pi}{6}\right]$)\\
		$k=1\quad \begin{cases}
			x_1 &=-\dfrac{\pi}{6}+2\pi=\dfrac{11\pi}{6}\in\left[\dfrac{0\pi}{3};\dfrac{6\pi}{3}\right]\\
			x_2 &=\dfrac{\pi}{6}+2\pi=\dfrac{13\pi}{6}\notin\left[\dfrac{0\pi}{3};\dfrac{6\pi}{3}\right]
		\end{cases}$\\
		On peut maintenant conclure : $S=\left\{\dfrac{\pi}{6};\dfrac{11\pi}{6}\right\}$
	\end{corrBox}
	\item $\sin(x)=-\dfrac{1}{2}$ avec $x\in \left[5\pi;7\pi\right]$\\
	\begin{corrBox}
		\begin{multicols}{2}
			Solution dans l'intervalle $]-\pi;\pi]$:\\
			$\begin{cases}
				x_1 &=-\dfrac{5\pi}{6}\\
				x_2 &=-\dfrac{\pi}{6}
			\end{cases}$\\
			
			\columnbreak
			On en déduit les solutions sur \R:\\
			$\begin{cases}
				x_1 &=-\dfrac{5\pi}{6}+2k\pi\\
				x_2 &=-\dfrac{\pi}{6}+2k\pi\
			\end{cases}$\\
		\end{multicols}
		On teste différentes valeurs de $k$ pour trouver les valeurs qui appartiennent à l'intervalle souhaité $\left[5\pi;7\pi\right]$ (ou $\left[\dfrac{30\pi}{6};\dfrac{42\pi}{6}\right]$)\\
		$k=3\quad \begin{cases}
			x_1 &=-\dfrac{5\pi}{6}+6\pi=\dfrac{31\pi}{6}\in\left[\dfrac{30\pi}{6};\dfrac{42\pi}{6}\right]\\
			x_2 &=-\dfrac{\pi}{6}+6\pi=\dfrac{35\pi}{6}\in\left[\dfrac{30\pi}{6};\dfrac{42\pi}{6}\right]
		\end{cases}$\\
		$k=4\quad \begin{cases}
			x_1 &=-\dfrac{5\pi}{6}+8\pi=\dfrac{43\pi}{6}\notin\left[\dfrac{30\pi}{6};\dfrac{42\pi}{6}\right]\\
			x_2 &=-\dfrac{\pi}{6}+8\pi=\dfrac{47\pi}{6}\notin\left[\dfrac{30\pi}{6};\dfrac{42\pi}{6}\right]
		\end{cases}$\\
		On peut maintenant conclure : $S=\left\{\dfrac{31\pi}{6};\dfrac{35\pi}{6}\right\}$
	\end{corrBox}
	\item $\cos(x)=\dfrac{\sqrt{2}}{2}$ avec $x\in \left[-3\pi;\pi\right]$\\
	\begin{corrBox}
		\begin{multicols}{2}
			Solution dans l'intervalle $]-\pi;\pi]$:\\
			$\begin{cases}
				x_1 &=-\dfrac{\pi}{4}\in\left[-\dfrac{12\pi}{4};\dfrac{4\pi}{4}\right]\\
				x_2 &=\dfrac{\pi}{4}\in\left[-\dfrac{12\pi}{4};\dfrac{4\pi}{4}\right]
			\end{cases}$\\
			
			\columnbreak
			On en déduit les solutions sur \R:\\
			$\begin{cases}
				x_1 &=-\dfrac{\pi}{4}+2k\pi\\
				x_2 &=\dfrac{\pi}{4}+2k\pi\
			\end{cases}$\\
		\end{multicols}
		On teste différentes valeurs de $k$ pour trouver les valeurs qui appartiennent à l'intervalle souhaité $\left[-3\pi;\pi\right]$ (ou $\left[-\dfrac{12\pi}{4};\dfrac{4\pi}{4}\right]$)\\
		$k=-1\quad \begin{cases}
			x_1 &=-\dfrac{\pi}{4}-2\pi=-\dfrac{9\pi}{4}\in\left[-\dfrac{12\pi}{4};\dfrac{4\pi}{4}\right]\\
			x_2 &=\dfrac{\pi}{4}-2\pi=-\dfrac{7\pi}{4}\in\left[-\dfrac{12\pi}{4};\dfrac{4\pi}{4}\right]
		\end{cases}$\\
		On peut maintenant conclure : $S=\left\{-\dfrac{9\pi}{4};-\dfrac{7\pi}{4};-\dfrac{\pi}{4};\dfrac{\pi}{4}\right\}$
	\end{corrBox}
	\item \ding{46} $\sin(x)=-\dfrac{\sqrt{2}}{2}$ avec $x\in \left[-8\pi;-6\pi\right]$\\
	\begin{corrBox}
		\begin{multicols}{2}
			Solution dans l'intervalle $]-\pi;\pi]$:\\
			$\begin{cases}
				x_1 &=-\dfrac{3\pi}{4}\\
				x_2 &=-\dfrac{\pi}{4}
			\end{cases}$\\
			
			\columnbreak
			On en déduit les solutions sur \R:\\
			$\begin{cases}
				x_1 &=-\dfrac{3\pi}{4}+2k\pi\\
				x_2 &=-\dfrac{\pi}{4}+2k\pi\
			\end{cases}$\\
		\end{multicols}
		On teste différentes valeurs de $k$ pour trouver les valeurs qui appartiennent à l'intervalle souhaité $\left[-8\pi;-6\pi\right]$ (ou $\left[-\dfrac{32\pi}{4};-\dfrac{24\pi}{4}\right]$)\\
		$k=-3\quad \begin{cases}
			x_1 &=-\dfrac{3\pi}{4}-6\pi=-\dfrac{27\pi}{4}\in\left[-\dfrac{32\pi}{4};-\dfrac{24\pi}{4}\right]\\
			x_2 &=-\dfrac{\pi}{4}-6\pi=-\dfrac{25\pi}{4}\in\left[-\dfrac{32\pi}{4};-\dfrac{24\pi}{4}\right]
		\end{cases}$\\
		$k=-4\quad \begin{cases}
			x_1 &=-\dfrac{3\pi}{4}-8\pi=-\dfrac{35\pi}{4}\notin\left[-\dfrac{32\pi}{4};-\dfrac{24\pi}{4}\right]\\
			x_2 &=-\dfrac{\pi}{4}-8\pi=-\dfrac{33\pi}{4}\notin\left[-\dfrac{32\pi}{4};-\dfrac{24\pi}{4}\right]
		\end{cases}$\\
		On peut maintenant conclure : $S=\left\{-\dfrac{27\pi}{4};-\dfrac{25\pi}{4}\right\}$
	\end{corrBox}
	\item $\cos(x)=\dfrac{17}{3}$ avec $x\in \left[-3\pi;7\pi\right]$\\
	\begin{corrBox}
		$\dfrac{17}{3}>1$ or pour tout $x \in \R$, $-1\leqslant \cos x\leqslant1$ donc\\
		$S=\varnothing$
	\end{corrBox}
\end{enumMet}
\end{document}

